% §13 Exercises
%
% Cross-reference resolution status:
%
% DEFINED here:
%   \label{ex:ch13-1} through \label{ex:ch13-11}
%
% RESOLVED (used here, defined elsewhere):
%   \Cref{sec:algebraic-types-recap}              -- §13.1 (Plan 13 Task 2)
%   \Cref{sec:bernoulli-of-composite}             -- §13.2 (Plan 13 Task 3)
%   \Cref{sec:composition-functor}                -- §13.3 (Plan 13 Task 4)
%   \Cref{thm:bernoulli-product}                  -- §13.2 (Plan 13 Task 3)
%   \Cref{thm:kronecker}                          -- ch4 §4.4
%   \Cref{thm:n-ary-kronecker}                    -- ch9 §9.3
%   \Cref{ch:approximate-functions}               -- ch9
%   \Cref{ch:sketches-estimators}                 -- ch10
%   \Cref{ch:approximate-relations}               -- ch12
%   \Cref{ch:not-bernoulli}                       -- ch16 (forward)
%
% Difficulty-marker macros (not in alex.sty; defined locally if not yet defined):
\providecommand{\exstar}{$\star$}
\providecommand{\exstarstar}{$\star\star$}
\providecommand{\exstarstarstar}{$\star\star\star$}

\section*{Exercises}
\addcontentsline{toc}{section}{Exercises}

\begin{enumerate}

% -----------------------------------------------------------------------
% WARM-UP  (★ × 3): exercises 1, 2, 3
% -----------------------------------------------------------------------

\item \exstar
\label{ex:ch13-1}
\textbf{Channel matrix for $\Bernoulli{\BoolType \times \BoolType}$.}
Let $\tilde b_1, \tilde b_2$ be two independent
$\Bernoulli{\BoolType}$ instances with channel matrices
\[
  Q_1 = \begin{pmatrix} 1 - \fnrate_1 & \fnrate_1 \\
                        \fprate_1 & 1 - \fprate_1 \end{pmatrix},
  \qquad
  Q_2 = \begin{pmatrix} 1 - \fnrate_2 & \fnrate_2 \\
                        \fprate_2 & 1 - \fprate_2 \end{pmatrix}.
\]
Each $Q_i$ follows the positives-first convention of
\Cref{def:bernoulli-bool}: row $1$ governs the positive input $\top$,
row $2$ the negative input $\bot$. The pair $(\tilde b_1, \tilde b_2)$
is a $\Bernoulli{\BoolType \times \BoolType}$ instance on the
$4$-element type, ordered positives-first to match:
$\{(\top, \top), (\top, \bot), (\bot, \top), (\bot, \bot)\}$.
Write the $4 \times 4$ joint channel matrix as the Kronecker product
$Q_1 \otimes Q_2$, with rows and columns indexed in the order above.
Verify by direct computation that the $(\top, \top) \to (\top, \top)$
entry equals $(1 - \fnrate_1)(1 - \fnrate_2)$, in agreement with
\Cref{thm:bernoulli-product}.

\item \exstar
\label{ex:ch13-2}
\textbf{Identify the type for each prior approximate structure.}
For each of the following, name the type $T$ such that the
structure is a $\Bernoulli{T}$ instance under the
\Cref{sec:algebraic-types-recap} taxonomy.
\begin{enumerate}
  \item[(a)] Bloom filter on universe $U$: a noisy membership
    predicate, hence $T = U \to \BoolType$.
  \item[(b)] CountMin sketch on key universe $K$: a noisy counter
    per key, hence $T = K \to \mathbb{N}$.
  \item[(c)] MinHash sketch on a set $A \subseteq U$: a noisy
    similarity estimate in $[0, 1]$ per pair of sets, hence
    $T = [0, 1]$ for the per-pair estimator (or
    $T = (U \to \BoolType) \times (U \to \BoolType) \to [0, 1]$ for
    the full pairwise-similarity table).
  \item[(d)] Approximate equijoin of $R_1 \subseteq T_1 \times T_2$
    and $R_2 \subseteq T_2 \times T_3$: a noisy relation on the
    composite pair universe, hence
    $T = (T_1 \times T_3) \to \BoolType$.
\end{enumerate}

\item \exstar
\label{ex:ch13-3}
\textbf{Block-diagonal channel for $\Bernoulli{\BoolType + \BoolType}$.}
The sum type $\BoolType + \BoolType$ has four elements:
$\mathrm{Left}(\bot)$, $\mathrm{Left}(\top)$, $\mathrm{Right}(\bot)$,
$\mathrm{Right}(\top)$. An instance of
$\Bernoulli{\BoolType + \BoolType}$ that preserves the tag (errors
within each summand but never confuses $\mathrm{Left}$ for
$\mathrm{Right}$) has a channel matrix of block-diagonal shape with
two $2 \times 2$ blocks along the diagonal and zeros off the
diagonal. Write the shape explicitly:
\[
  Q = \begin{pmatrix} Q_L & 0 \\ 0 & Q_R \end{pmatrix},
  \qquad Q_L, Q_R \in \mathbb{R}^{2 \times 2},
\]
where $Q_L$ and $Q_R$ are the per-summand $\Bernoulli{\BoolType}$
channels. State in one sentence the structural meaning: a
tag-preserving noisy sum is two independent noisy bits selected by
the tag.

% -----------------------------------------------------------------------
% MEDIUM  (★★ × 5): exercises 4, 5, 6, 7, 8
% -----------------------------------------------------------------------

\item \exstarstar
\label{ex:ch13-4}
\textbf{Composition isomorphism for products.}
Prove $\Bernoulli{T \times U} \cong \Bernoulli{T} \otimes
\Bernoulli{U}$ at the channel-matrix level, using
\Cref{thm:kronecker} and \Cref{thm:n-ary-kronecker}.
\begin{enumerate}
  \item[(a)] Forward direction: given a $\Bernoulli{T}$ with channel
    matrix $Q_T$ and a $\Bernoulli{U}$ with channel matrix $Q_U$,
    and given that the two are independent, show that the joint
    channel matrix on $T \times U$ is exactly $Q_T \otimes Q_U$
    by writing out the entry-wise product
    $\Pr[\hat t, \hat u \mid t, u] = \Pr[\hat t \mid t]
    \cdot \Pr[\hat u \mid u]$ and matching it to the Kronecker
    definition.
  \item[(b)] Backward direction: given any rank-$1$-factorizable
    channel matrix $Q$ on $T \times U$ (one whose entries factor as
    $Q_{(j_T, j_U), (\ell_T, \ell_U)} = a_{j_T, \ell_T} \cdot
    b_{j_U, \ell_U}$), recover the marginal channels $Q_T$ and $Q_U$
    as $a$ and $b$ respectively, normalised so each row sums to $1$.
  \item[(c)] Conclude that the isomorphism class of
    $\Bernoulli{T \times U}$ instances under independent factors is
    in bijection with pairs of $\Bernoulli{T}$ and $\Bernoulli{U}$
    instances. The non-independent case (which lives outside the
    bijection) is the subject of \Cref{ch:not-bernoulli}.
\end{enumerate}

\item \exstarstar
\label{ex:ch13-5}
\textbf{Sketches as $\Bernoulli{T \to \BoolType}$ vs $\Bernoulli{T}$.}
The sketches of \Cref{ch:sketches-estimators} divide into two type
classes by codomain.
\begin{enumerate}
  \item[(a)] Bloom filters and quotient filters are
    $\Bernoulli{T \to \BoolType}$ instances: noisy predicates,
    one Boolean per key. Identify the per-key Boolean as the
    membership indicator and the noise as the
    $\Bernoulli{\BoolType}$ on it.
  \item[(b)] CountMin and HyperLogLog are $\Bernoulli{T \to
    \mathbb{N}}$ instances: noisy counters or cardinality estimators
    per key, one natural number per key. Identify the
    per-key $\mathbb{N}$-valued attribute and the noise as the
    $\Bernoulli{\mathbb{N}}$ on it.
  \item[(c)] State which row of the
    \Cref{sec:composition-functor} table determines composition
    for each class: the function row $\Bernoulli{T \to U} \cong
    T \to \Bernoulli{U}$ in both cases, with $U = \BoolType$ for
    (a) and $U = \mathbb{N}$ for (b). The composition rules
    differ only in the per-output $\Bernoulli{U}$ algebra, not in
    the per-input type structure.
\end{enumerate}

\item \exstarstar
\label{ex:ch13-6}
\textbf{$\Bernoulli{-}$ as a monad.}
Explain informally why the isomorphism $\Bernoulli{T \to U} \cong
T \to \Bernoulli{U}$ from \Cref{sec:composition-functor} qualifies
$\Bernoulli{-}$ as a ``monad'' in the programming-language sense.
The argument should cover:
\begin{enumerate}
  \item[(a)] The Kleisli-arrow shape: a noise-introducing function
    has signature $T \to \Bernoulli{U}$, with input a pure value and
    output a noisy value of a possibly different type. This is the
    monad's $\Rightarrow$ arrow.
  \item[(b)] Composition of Kleisli arrows: two such arrows
    $f : T \to \Bernoulli{U}$ and $g : U \to \Bernoulli{V}$ compose
    into $g \circ f : T \to \Bernoulli{V}$ by feeding the
    $\Bernoulli{U}$ output of $f$ through $g$ pointwise and
    averaging out the intermediate noise. The averaging step is the
    monad's join (or bind).
  \item[(c)] Unit: every pure value $t : T$ embeds into
    $\Bernoulli{T}$ as the point distribution concentrated on $t$
    (the noiseless special case of a $\Bernoulli{T}$ instance). This
    is the monad's unit.
\end{enumerate}
The three pieces (Kleisli arrow, composition by averaging, unit by
point mass) make $\Bernoulli{-}$ a monad in the
programming-language sense; the categorical axioms are
in \textcite{wadler-1992}.

\item \exstarstar
\label{ex:ch13-7}
\textbf{Composition rules for lists.}
Consider $\Bernoulli{L T}$ where the list length is itself noisy
(as in a HyperLogLog-style cardinality estimate combined with
per-element approximate values). Derive the composition rule.
\begin{enumerate}
  \item[(a)] Factor the type: $L T \cong \sum_{n \ge 0} T^n$, the
    sum over $n$ of fixed-length $n$-tuples. Apply the sum row of
    the \Cref{sec:composition-functor} table to factor
    $\Bernoulli{L T}$ as a block-diagonal channel indexed by length,
    with the $n$-th block a $\Bernoulli{T^n}$.
  \item[(b)] Apply the product row to each block: each
    $\Bernoulli{T^n}$ factors as $\Bernoulli{T}^{\otimes n}$ under
    per-position independence.
  \item[(c)] Combine: $\Bernoulli{L T}$ is a noisy length (a
    $\Bernoulli{\mathbb{N}}$) plus, conditional on length $n$, a
    Kronecker product of $n$ independent $\Bernoulli{T}$ instances.
    The full composition rule is a mixture over the noisy length,
    with each mixture component a per-position lift. State in one
    sentence the implication for HyperLogLog: the length-noise
    (cardinality estimate) and the per-element noise compose by
    this rule, not by a simple per-element product.
\end{enumerate}

\item \exstarstar
\label{ex:ch13-8}
\textbf{Type for an approximate database.}
A database has tables (each a set of tuples), indexes (each a map
from a key attribute to a tuple identifier), and queries (each a
function from a query specification to a result set). Identify the
type $T$ for which an approximate database is a $\Bernoulli{T}$
instance, by decomposing the database structure into the four type
constructors of \Cref{sec:algebraic-types-recap}.
\begin{enumerate}
  \item[(a)] A single table is a list of tuples, so $T_{\text{table}}
    = L (A_1 \times A_2 \times \cdots \times A_k)$ for attribute
    types $A_1, \ldots, A_k$.
  \item[(b)] A single index is a function from a key value to a
    tuple identifier (or absence), so
    $T_{\text{index}} = \text{Key} \to (\text{TupleId} + \text{Unit})$.
  \item[(c)] The database as a whole is a product of named tables
    and indexes, plus a query function: $T_{\text{db}} =
    T_{\text{table}}^{(1)} \times \cdots \times T_{\text{table}}^{(m)}
    \times T_{\text{index}}^{(1)} \times \cdots \times
    T_{\text{index}}^{(p)} \times (\text{QuerySpec} \to L
    \text{Tuple})$. An approximate database is a $\Bernoulli{}$ on
    this composite type, with each constructor's noise governed by
    its row of the \Cref{sec:composition-functor} table.
\end{enumerate}

% -----------------------------------------------------------------------
% CHALLENGING / EXTENSION  (★★★ × 3): exercises 9, 10, 11
% -----------------------------------------------------------------------

\item \exstarstarstar
\label{ex:ch13-9}
\textbf{Categorical functor.}
State the technical definition of a \emph{functor} from category
theory, and verify that $\Bernoulli{-}$ satisfies it for the four
type constructors.
\begin{enumerate}
  \item[(a)] Definition: a functor $F : \mathcal{C} \to \mathcal{D}$
    assigns to each object $X$ of $\mathcal{C}$ an object $F(X)$ of
    $\mathcal{D}$, and to each morphism $f : X \to Y$ a morphism
    $F(f) : F(X) \to F(Y)$, such that $F(\mathrm{id}_X) =
    \mathrm{id}_{F(X)}$ and $F(g \circ f) = F(g) \circ F(f)$ for
    composable $f, g$.
  \item[(b)] Take $\mathcal{C}$ to be the category of algebraic
    types with type morphisms (functions $T \to U$) and
    $\mathcal{D}$ to be the category of noisy channels. The action
    on objects is $T \mapsto \Bernoulli{T}$. The action on
    morphisms is $(f : T \to U) \mapsto (\Bernoulli{f} :
    \Bernoulli{T} \to \Bernoulli{U})$, where $\Bernoulli{f}$
    transports a $\Bernoulli{T}$ instance through $f$ pointwise
    (pushforward of the channel matrix).
  \item[(c)] Verify the two functor axioms: identity preservation
    ($\Bernoulli{\mathrm{id}_T} = \mathrm{id}_{\Bernoulli{T}}$ since
    the identity function induces the identity channel) and
    composition preservation ($\Bernoulli{g \circ f} = \Bernoulli{g}
    \circ \Bernoulli{f}$ since channel pushforward is functorial in
    the underlying function). The four type-constructor identities
    of \Cref{sec:composition-functor} then say that $\Bernoulli{-}$
    is moreover a monoidal functor with respect to each of the four
    constructors.
\end{enumerate}

\item \exstarstarstar
\label{ex:ch13-10}
\textbf{Universal property of $\Bernoulli{-}$.}
State a universal property characterizing $\Bernoulli{-}$ up to
isomorphism, in the spirit of free constructions in algebra (free
monoid, free abelian group, free monad on a signature).
\emph{One approach:} $\Bernoulli{-}$ is the free monad on the
signature of a single binary symmetric channel.
\begin{enumerate}
  \item[(a)] The signature has one operation, the binary symmetric
    channel $\mathrm{bsc} : \BoolType \to \Bernoulli{\BoolType}$
    of Part~I, with two parameters $(\fprate, \fnrate)$.
  \item[(b)] The free monad on this signature is the monad whose
    Kleisli arrows are precisely the compositions of $\mathrm{bsc}$
    instances (with parameters varying) and pure functions. Argue
    that this free monad is isomorphic to $\Bernoulli{-}$ on the
    subcategory of types built from $\BoolType$ and the four
    constructors.
  \item[(c)] The universal property is then the standard
    free-construction one: any monad $M$ on the same category
    equipped with a binary-symmetric-channel operation of its own
    factors uniquely through $\Bernoulli{-}$. State the property
    precisely (the unique factorization condition) and identify
    why $\Bernoulli{-}$ is the most economical such monad: it adds
    no structure beyond what the channel signature forces.
\end{enumerate}

\item \exstarstarstar
\label{ex:ch13-11}
\textbf{Boundary case for the algebra.}
Construct an approximate structure that does \emph{not} fit any of
the four type constructors of \Cref{sec:algebraic-types-recap}, and
hence falls outside the algebra of \Cref{sec:composition-functor}.
\Cref{ch:not-bernoulli} catalogues such cases; the exercise asks
for one explicit construction.
\begin{enumerate}
  \item[(a)] Consider an approximate structure with correlated noise
    across positions: a noisy list $(\tilde t_1, \ldots, \tilde t_n)$
    whose entries are not independent. The list row of the
    \Cref{sec:composition-functor} table assumes per-position
    independence, so a correlated noisy list does not factor as
    $L \Bernoulli{T}$.
  \item[(b)] Alternatively: a structure with a noisy type tag
    (uncertainty over which type constructor applies). For example,
    a query result that is either ``a $\Bernoulli{\text{Hit}}$ with
    payload'' or ``a $\Bernoulli{\text{Miss}}$ with no payload,''
    where the tag itself is noisy. The sum row of the table assumes
    a deterministic tag, so a noisy-tag sum does not factor as
    $\Bernoulli{T} \oplus \Bernoulli{U}$.
  \item[(c)] State in one sentence why these cases require
    \Cref{ch:not-bernoulli}: they introduce dependence structure
    that no factorization of the channel matrix into type-level
    components can express, so the four-row algebra is incomplete.
\end{enumerate}

\end{enumerate}
